\SGAfrontchapter{原书题名与版本说明}
\SourcePage{1}
\begin{center}
{\Large 布瓦玛丽代数几何讨论班}\par
1960—1961\par\bigskip
{\LARGE\bfseries \Term{etale-covering}{平展覆盖}与\Term{fundamental-group}{基本群}}\par
（SGA 1）\par\bigskip
由 A. Grothendieck 主持的讨论班\par
增补 M. Raynaud 女士撰写的两讲\par\bigskip
本书为 Springer-Verlag 于 1971 年出版的\par
《Lecture Notes in Mathematics》第 224 卷的重排注释版。\par
原文件版本标识：arXiv:math/0206203v2，2004 年 1 月 4 日。
\end{center}

\SGAfrontchapter{重排版序言}
\SourcePage{3}
本文是 Alexander Grothendieck 等人所著《Revêtements Étales et Groupe Fondamental》（《\Term{etale-covering}{平展覆盖}与\Term{fundamental-group}{基本群}》）的重排注释版。原书为《Lecture Notes in Mathematics》第 224 卷，1971 年由 Springer-Verlag 在柏林、海德堡、纽约出版。

\LaTeXe{} 排版工作由参加 Bas Edixhoven 所主持项目的志愿者完成；该项目的进一步信息可见 \url{http://www.math.leidenuniv.nl/~edix/}。最后的版面整理由法国数学会完成。

本版对原文作了少量订正，由法国数学会在“Documents Mathématiques”丛书中出版（\url{http://smf.emath.fr/}）。Michel Raynaud 补充了若干更新说明；这些说明用方括号括起，并标以（MR），见第 211、220、256、292 页的注记及第 67 页注 7。为统一记号，点 $x$ 的\Term{residue-field}{剩余域}记为 $k(x)$，\Term{local-ring}{局部环} $A$ 的\Term{residue-field}{剩余域}记为 $k(A)$。

另有一种电子版本，旨在重现原始文本。两个版本共用同一个源文件，存放于 arXiv.org 电子预印本服务器 \url{http://arxiv.org/}；两个版本之间的差异均在源文件中说明。

原版页码标在页边，数字 $n$ 表示原版第 $n$ 页的起始位置。

\SGAfrontchapter{重排版英文序言（译文）}
\SourcePage{4}
本文是 Alexander Grothendieck 等人所著《Revêtements Étales et Groupe Fondamental》的更新版。原书为《Lecture Notes in Mathematics》第 224 卷，1971 年由 Springer-Verlag 在柏林、海德堡、纽约出版。

\LaTeXe{} 排版由参加 Bas Edixhoven 所主持项目的志愿者完成，项目信息见 \url{http://www.math.leidenuniv.nl/~edix/}；法国数学会完成了后续工作。

本版对原文作了少量订正，由法国数学会作为“Documents Mathématiques”丛书的一卷出版。Michel Raynaud 增加了若干更新说明；这些说明用方括号括起，并标以（MR），见第 211、220、256、292 页的注记及第 67 页注 7。为统一记号，点 $x$ 的\Term{residue-field}{剩余域}记为 $k(x)$，\Term{local-ring}{局部环} $A$ 的\Term{residue-field}{剩余域}记为 $k(A)$。

另有一个旨在重现原始文本的电子版本。这两个版本由同一个源文件生成；该文件可从 arXiv.org 电子预印本服务器 \url{http://arxiv.org/} 下载。两版之间的全部差异均在源文件中说明。

原版页码标在页边，数字 $n$ 表示原版第 $n$ 页的起始位置。

\SGAfrontchapter{引言}
\SourcePage{5}
本引言第一部分说明本卷的内容，第二部分介绍“布瓦玛丽代数几何讨论班”全套丛书；本卷是其中的第一卷。

\textbf{1.} 本卷从“克罗内克式”的观点出发，奠定代数几何中\Term{fundamental-group}{基本群}理论的基础。这一观点使我们能够同样处理通常意义下的代数簇，以及数域的整数环等对象。只有使用\Term{scheme}{概形}语言，才能令人满意地表述这一观点。因此，我们将自由使用这一语言，以及 J. Dieudonné 与 A. Grothendieck 的《Éléments de Géométrie Algébrique》（《代数几何基础》，下文简称 EGA）前三章所陈述的主要结果。阅读本卷不需要其他代数几何知识；对于希望熟悉当代代数几何技术的读者，本卷也可作为这些技术的入门。

本书第 I—XI 讲几乎原封不动地重印了口头讨论班的油印讲义；这些讲义原由高等科学研究所负责分发。\footnote{此后各次讨论班的讲义亦如此分发。这种分发方式逐渐显得难以实行且不足以满足需要，因而今后“布瓦玛丽代数几何讨论班”的各部分都将像本卷一样以书籍形式出版。}我们只在原文中增添了少量脚注，订正了一些打字错误，并调整了术语；特别是将其间已弃用的“简单态射”（morphisme simple）替换为“\Term{smooth-morphism}{光滑态射}”（morphisme lisse），以免发生同样的混淆。

\SourcePage{6}
第 I—IV 讲介绍\Term{etale-morphism}{平展态射}和\Term{smooth-morphism}{光滑态射}这两个局部概念，几乎不超出《基础》第一章所阐述的\Term{scheme}{概形}语言。\footnote{目前《基础》第 IV 章第 17、18 节已有更全面的论述。}第 V 讲给出\Term{scheme}{概形}的\Term{fundamental-group}{基本群}的公理化描述；即使在\Term{scheme}{概形}退化为一个域的谱这一经典情形下，这一描述仍很有用，并给出通常 Galois 理论的一种十分方便的重新表述。第 VI、VIII 讲介绍\Term{descent}{下降}理论；近年来，这一理论在代数几何中的重要性日益增加，在解析几何和拓扑学中也可能发挥类似的作用。应当指出，第 VII 讲没有写成讲义，其内容已经收入 J. Giraud 的著作《Méthode de la Descente》（《法国数学会公报》，专刊第 2 号，1964 年，viii + 150 页）。第 IX 讲专门研究\Term{etale-morphism}{平展态射}的\Term{descent}{下降}，由此系统处理\Term{fundamental-group}{基本群}的 Van Kampen 型定理；在这里，这些定理只是\Term{descent}{下降}定理的另一种表述。具体而言，设\Term{connected}{连通}\Term{scheme}{概形} $X$ 上给定一个\Term{surjectivity}{满射}且\Term{proper}{固有}的\Term{morphism}{态射} $X'\to X$，我们可以借助 $X'$ 及\Term{fiber-product}{纤维积}
\[
X'\times_X X',\qquad X'\times_X X'\times_X X'
\]
的各个\Term{connected-component}{连通分支}的\Term{fundamental-group}{基本群}，以及上述\Term{scheme}{概形}之间的典范单纯\Term{morphism}{态射}所诱导的群\Term{homomorphism}{同态}，来计算 $X$ 的\Term{fundamental-group}{基本群}。第 X 讲给出\Term{proper}{固有}\Term{smooth-morphism}{光滑态射}下\Term{fundamental-group}{基本群}的\Term{specialization}{特殊化}理论；其中最突出的结果是：借助特征零时用超越方法所得的已知结论，几乎完全确定特征 $p>0$ 时\Term{smooth}{光滑}代数曲线的\Term{fundamental-group}{基本群}。第 XI 讲给出一些例子和补充，特别是以\Term{cohomology}{上同调}形式明确表述 Kummer \Term{covering}{覆盖}理论和 Artin–Schreier \Term{covering}{覆盖}理论。关于正文的其他说明，见本引言之后的《油印版说明》。

自本讨论班于 1961 年写成以来，我与 M. Artin 合作发展了\Term{etale-topology}{平展拓扑}的语言及相应的\Term{cohomology}{上同调}理论。这些内容载于代数几何讨论班的 SGA 4 部分《\Term{scheme}{概形}的\Term{etale-cohomology}{平展上同调}》，即将在与本卷相同的丛书中出版。这套语言及目前已有的结果，为研究\Term{fundamental-group}{基本群}提供了格外灵活的工具，使人能够更好地理解并超越这里陈述的部分结果。因此，\Term{fundamental-group}{基本群}理论应当从这一观点出发重新作一次全面阐述；事实上，所有关键结果如今都已见于上述文献。原拟在《基础》中专门讨论\Term{fundamental-group}{基本群}的一章正要完成这项工作，并收入其他一些这里未能容纳的发展，这些发展依赖于消解奇点的技术：通过\Term{local-ring}{局部环}的适当\Term{resolution-singularities}{奇点消解}，计算完备\Term{local-ring}{局部环}的“\Term{local-fundamental-group}{局部基本群}”；在没有\Term{proper}{固有}性假设时给出\Term{fundamental-group}{基本群}的局部和整体 Künneth 公式（参见第 XIII 讲）；以及 M. Artin 关于\Term{excellent}{优秀} Hensel \Term{local-ring}{局部环}与其\Term{completion}{完备化}的\Term{local-fundamental-group}{局部基本群}之比较的结果（SGA 4 XIX）。

\SourcePage{7}
还应指出，有必要发展\Term{topos}{拓扑斯}的\Term{fundamental-group}{基本群}理论，将通常的拓扑理论、它的半单纯版本、本卷第 V 讲发展的“\Term{profinite}{投有限}”版本，以及 SGA 3 X 7 中稍为一般的投\Term{discrete}{离散}版本统摄起来；最后这个版本适用于非\Term{normal}{正规}且非单分支\Term{scheme}{概形}。在依此观点全面重写理论之前，Raynaud 女士的第 XIII 讲利用 SGA 4 的语言和结果，展示了在一些典型问题中这些工具的作用，特别是将第 X 讲的某些结果推广到非\Term{proper}{固有}的相对\Term{scheme}{概形}。该讲尤其给出了任意特征下非完备代数曲线的“素于 $p$”\Term{fundamental-group}{基本群}的结构；我在 1959 年曾宣布这一结果，但迄今尚未发表证明。

尽管存在这些缺漏和不足——也有人会说，恰恰因为存在这些缺漏和不足——我仍认为，在一部能够免受上述批评的著作写成并出版之前，本卷既可帮助希望熟悉\Term{fundamental-group}{基本群}理论的读者，也可充当参考书。

\textbf{2.} 本卷是“布瓦玛丽代数几何讨论班”的第一卷；后续各卷拟在同一丛书中出版。本讨论班与 J. Dieudonné、A. Grothendieck 的《代数几何基础》并行，其目标是按照后者所持的观点建立代数几何的基础。本讨论班各卷的标准参考文献是《基础》第 I、II、III 章（记为 EGA I、II、III），并假定读者掌握这些章节所需的交换代数和同调代数知识。\footnote{对此的详细说明见 EGA I 的引言。}此外，各卷将按需要自由引用本讨论班此前各卷，以及《基础》中已出版或即将出版的其他章节。

讨论班的每一部分都围绕一个主要主题展开，相应一卷或数卷的书名即表明这一主题。口头讨论班通常持续一个学年，有时更长。同一部分内部的各讲一般在逻辑上紧密相依；不同部分之间则在很大程度上相互独立。例如，“\Term{group-scheme}{群概形}”部分几乎完全独立于此前的两部分，却经常使用 EGA IV 的结果。下面列出即将出版的各部分；下文记为 SGA 1 至 SGA 7：
\SourcePage{8}
\begin{description}[style=nextline,leftmargin=1.8em]
\item[SGA 1.] \Term{etale-covering}{平展覆盖}与\Term{fundamental-group}{基本群}，1960、1961 年。
\item[SGA 2.] \Term{coherent}{凝聚}层的局部\Term{cohomology}{上同调}与局部及整体 Lefschetz 定理，1961—1962 年。
\item[SGA 3.] \Term{group-scheme}{群概形}，1963、1964 年（三卷），与 M. Demazure 合作。
\item[SGA 4.] \Term{topos}{拓扑斯}理论与\Term{scheme}{概形}的\Term{etale-cohomology}{平展上同调}，1963—1964 年（三卷），与 M. Artin、J. L. Verdier 合作。
\item[SGA 5.] $\ell$ 进\Term{cohomology}{上同调}与 $L$ 函数，1964、1965 年（两卷）。
\item[SGA 6.] 相交理论与 Riemann–Roch 定理，1966—1967 年（两卷），与 P. Berthelot、L. Illusie 合作。
\item[SGA 7.] 代数几何中的局部单值群。
\end{description}
上述部分中有三部是与其他数学家共同主持的，这些合作者将列为相应各卷封面的共同作者。至于讨论班其他积极参与者，他们在讲义撰写和数学内容的完善两方面所起的作用逐年增加；每位参与者的姓名列于其作为报告人或撰稿人所负责的讲义开头，一卷所涉及的人员名单则列于该卷的扉页。

关于讨论班与《基础》的关系，有必要再作一些说明。《基础》原则上拟将代数几何中被认为最根本的概念与技术作一系统阐述；随着逻辑协调与美学要求的自然作用，这些概念和技术本身也逐渐明确。从这一观点出发，很自然会把讨论班视为《基础》的预备版本，迟早将几乎全部收入后者。几年前，这个过程已在一定程度上开始：本卷 SGA 1 的第 I—IV 讲已全部包含在 EGA IV 中，而第 VI—VIII 讲原拟在数年内收入 EGA VI。然而，随着《基础》和讨论班中所开展的建设工作不断推进，整体规模逐渐显现，最初把讨论班仅视作预备性、临时性版本的原则越来越显得不切实际。原因之一，是深谋远虑的大自然对人类寿命施加了限制。鉴于讨论班各部分在写作时通常已相当用心，今后只有在写作后的进一步进展足以通过较深入的改动带来极其显著的改善时，才有必要在《基础》或接替它的其他论著中重新阐述。

\SourcePage{9}
如上所述，本讨论班 SGA 1 目前已属这种情况；SGA 2 也是如此，这要归功于 Raynaud 女士的近期结果。反之，目前没有迹象表明，在不远的将来，前述 SGA 3 至 SGA 7 中任何一部分也会出现这种情况。

“布瓦玛丽代数几何讨论班”内部的引文约定如下。引用 SGA 1 至 SGA 7 某一部分内部的结果，采用 III 9.7 这样的形式；罗马数字 III 表示讲次，见该讲各页页眉，9.7 则表示该讲内部命题、定义、注记等的编号。必要时使用更长的小数式编号，例如 9.7.1、9.7.2，\Term{representable}{可表}示命题 9.7 证明中的各个步骤。III 9 表示第 III 讲第 9 节。同一讲内部的引用略去讲次。引用讨论班其他部分时，采用同样格式，但在前面注明相应部分，例如 SGA 1 III 9.7。同样，EGA IV 11.5.7 表示《代数几何基础》第 IV 章的命题、定义等 11.5.7；其中第一个阿拉伯数字仍表示节号。除这些全套 SGA 通用的约定之外，每讲涉及的参考文献通常集中列在该讲末尾，讲内按惯例以方括号内的数字引用，例如 [3]。

最后，为方便读者，在认为有必要时，我们将在各卷末尾附记号索引和术语索引，并酌情在术语索引中给出所用法文术语的英文译名。

我还想在这篇引言中补充一段数学以外的话。1969 年 11 月，我得知高等科学研究所——我基本上从它成立之初起就在其中任教授——已连续三年接受军事部的资助。我刚开始从事研究时，就已对大多数科学家毫不介意地以各种形式与军事机构合作感到极为遗憾。那时我的动机主要出于道德，因而很难受到认真对待。如今，军事机构及其所掌握的大规模毁灭手段不断\Term{group-extension}{扩张}，威胁着人类的存续，使这些动机具有了新的力量与意义。我已在别处更详细地谈过这些问题；在我看来，它们远比任何科学——包括数学——的进展更为重要。例如，可参阅《Survivre》杂志第 1 期（1970 年 8 月）G. Edwards 的文章，其中概述了我在别处就这些问题所作的较详细陈述。就这样，我在一个机构工作了三年，却不知道它采用了一种我认为不道德而且危险的资助方式。\footnote{不言而喻，我在这里表达的意见仅由我本人负责，与出版本卷的 Springer 出版社无关。}

\SourcePage{10}
目前，在 IHES 的同事中，只有我持这种意见，因此我为取消 IHES 预算中的军事资助所作的努力没有成功。于是，我作出了我认为应当作出的决定：于 1970 年 9 月 30 日离开 IHES，并在该机构继续接受这类资助期间，中止与它的一切科学合作。我已请求 IHES 所长 Motchane 先生，自 1970 年 10 月 1 日起，IHES 不再分发由我撰写的数学文献，或属于“布瓦玛丽代数几何讨论班”的文献。如前所述，本讨论班将由 Julius Springer 出版社在 Lecture Notes 丛书中出版。我很高兴在此感谢 Springer 出版社及 K. Peters 先生；他们高效而友善的帮助使这一出版工作得以实现，特别是承担了增补到旧讨论班中的新讲义，以及原来未齐全的讨论班中缺失讲义的胶印打字工作。

我也感谢 J. P. Delale 先生，他承担了编制记号索引和术语索引这项费力的工作。

\begin{flushright}马西，1970 年 8 月。\end{flushright}

\SGAfrontchapter{油印版说明}
\SourcePage{11}
每篇已写成的讲义都包含若干次连续口头报告的实质内容，因而没有必要逐一注明日期。

第 VIII 讲多次引用的第 VII 讲没有由报告人写成讲义。口头报告仅从严格实用的角度，概述了一般\Term{category}{范畴}中\Term{descent}{下降}的语言，没有讨论这种语言所引起的逻辑困难。后来发现，即使只考虑篇幅，对这一语言作出正确阐述也会超出本讲义的范围。关于\Term{descent}{下降}理论的正式论述，请参阅 Jean Giraud 正在准备的文章。在该文发表之前，\footnote{现已出版：J. Giraud，\textit{Méthodes de la descente}，法国数学会专刊第 2 号，1964 年。}我相信细心的读者能够自行补足第 VIII 讲中这些并无相应正文的引用。

第 XI 讲之后的另一些口头报告也没有写成讲义，正文有些地方曾提到它们；这些报告原拟构成第 XII、XIII 讲。前面几次报告在允许\Term{nilpotent}{幂零}元的\Term{scheme}{概形}和\Term{analytic-space}{解析空间}的框架中——这种\Term{analytic-space}{解析空间}是在 1960—1961 年 Cartan 讨论班中引入的——重新讨论了以下内容：对完备赋值域 $k$ 上\Term{locally-finite-type}{局部有限型}的\Term{prescheme}{预概形}，构造其\Term{associated-analytic-space}{相伴解析空间}；在 $k$ 为复数域时建立 GAGA 型定理；以及应用这些结果，比较用超越方法定义的\Term{fundamental-group}{基本群}与本讲义所研究的\Term{fundamental-group}{基本群}。参见 A. Grothendieck，\textit{Fondements de la Géométrie Algébrique}，Bourbaki 讨论班第 190 讲，第 10 页，1959 年 12 月。后面几次口头报告则概述了如何推广正文中发展的方法，以研究允许\Term{tame-ramification}{温和分歧}的\Term{covering}{覆盖}，以及从完备曲线删去\Term{finite}{有限}多个点后所得曲线的\Term{fundamental-group}{基本群}结构；参见上述文献第 182 讲，第 27 页，定理 14。这些报告并没有引入本质上新的思想，因此，在《代数几何基础》相应各章问世之前，似乎没有必要将它们正式写出。\footnote{这些内容已收入本卷 Raynaud 女士撰写的第 XII 讲，其中的证明不同于口头讨论班中的原始证明，参见引言。}

\SourcePage{12}
另一方面，\Term{fundamental-group}{基本群}和\Term{picard-group}{Picard 群}的局部及整体 Lefschetz 型定理，已在 1962 年的另一届讨论班中专门讨论，且已全部写成讲义，可供读者使用。\footnote{《Cohomologie étale des faisceaux cohérents et théorèmes de Lefschetz locaux et globaux》（SGA 2），由 North Holland Pub. Cie. 出版。〔译注：此处题名中的 étale 为原文所载；本套 SGA 2 的题名使用 locale。〕}需要指出，本讨论班和 1962 年讨论班中发展的结果，将在证明\Term{prescheme}{预概形}\Term{etale-cohomology}{平展上同调}中的若干关键结论时发挥根本作用；这些结论将构成由 M. Artin 与我共同主持的 1963—1964 年讨论班的内容，该讨论班目前正在筹备中。\footnote{《Cohomologie étale des schémas》（记为 SGA 4），即将在本丛书出版。}

第 I—IV 讲本质上是局部的，且非常初等，它们将全部收入《代数几何基础》第 IV 章。该章第一部分正在印刷，预计于 1964 年底出版。尽管如此，对于希望先了解\Term{smooth-morphism}{光滑态射}、\Term{etale-morphism}{平展态射}或\Term{flat-morphism}{平坦态射}的基本性质，再深入系统论著的读者，这几讲仍可能有用。至于其余各讲，则将收入《基础》第 VIII 章；\footnote{事实上，由于《基础》原定计划发生了变化，\Term{fundamental-group}{基本群}的研究已推迟到比此处所指更后的一章。参见本《说明》前的引言。}该章至少还需数年才有可能出版。

\begin{flushright}比尔，1963 年 6 月。\end{flushright}
